\section{Step 2: Identify the analytic quantity that can have a resonance}
\label{sec:resonance-definition}
% BEGIN CURATED SUBJECT INDEX
\index{spectrum!point}
\index{resolvent}
\index{Jost function}
\index{scattering matrix}
\index{boundary condition!incoming and outgoing}
\index{resonance!residue}
% END CURATED SUBJECT INDEX

\calculationroute
  {Jost solutions, resolvent boundary values, and the physical-sheet
  scattering matrix.}
  {Continue the incoming coefficient through the continuum cut and relate its
  zero to resolvent and $S$-matrix poles.}
  {State the sheet, boundary condition, and residue associated with a proposed
  resonance rather than identifying it only from a complex eigenvalue.}

\subsection{Self-adjoint starting point and resolvent boundary values}
\label{subsec:resolvent}
% BEGIN CURATED SUBJECT INDEX
\index{operator!adjoint and self-adjoint}
\index{Hilbert space}
\index{topology}
\index{spectrum!continuous}
\index{resolvent}
\index{GNS construction}
\index{scattering matrix}
\index{boundary condition!incoming and outgoing}
\index{Riemann sheet}
\index{resonance!residue}
% END CURATED SUBJECT INDEX

Let $\Hil$ be a separable Hilbert space and let $H$ be a self-adjoint
Hamiltonian on a dense domain $D(H)\subset\Hil$.  For a complex number
$z\notin\spectrum(H)$, where $\spectrum(H)$ is the spectrum, define the resolvent
\begin{equation}
  \Resolvent(z)=(z-H)^{-1} .
  \label{eq:resolvent-definition}
\end{equation}
The operator-valued resolvent of a self-adjoint $H$ is analytic away from the
real spectrum.  If $H$ has continuous spectrum on an interval $I\subset
\real$, the limits $\Resolvent(E\pm\rmi0)$ are generally meaningful only in a
weighted space or as matrix elements between sufficiently localized test
vectors.  The two signs are the outgoing and incoming boundary values.

Choose test vectors $\phi,\psi$ from a dense space $\PhiSpace$ with stronger
topology than $\Hil$.  A resonance is a pole $z=\vsub{z}{R}$ of a
meromorphic continuation of
\begin{equation}
  F_{\phi\psi}(z)
  =\bra{\phi}\Resolvent(z)\ket{\psi}
  \label{eq:resolvent-matrix-element}
\end{equation}
through the continuous spectrum to an unphysical sheet.  For a simple pole,
\begin{equation}
  F_{\phi\psi}(z)
  =\frac{\dualpair{\phi}{\vsub{\Psi}{R}}
  \dualpair{\vsub{\widetilde{\Psi}}{R}}{\psi}}
  {z-\vsub{z}{R}}
  +\vsup{F_{\phi\psi}}{reg}(z) .
  \label{eq:resolvent-laurent}
\end{equation}
The right and left pole functionals are denoted by
$\vsub{\Psi}{R}$ and $\vsub{\widetilde{\Psi}}{R}$, respectively, and
$\vsup{F_{\phi\psi}}{reg}$ is regular at $\vsub{z}{R}$.  For a
time-reversal-invariant one-channel problem, the pole energy is conventionally
written
\begin{equation}
  \vsub{z}{R}=E_*-\frac{\rmi}{2}\Gamma,
  \qquad E_*\in\real,
  \qquad \Gamma>0 .
  \label{eq:pole-energy}
\end{equation}
The pole lifetime is $\tau=\hbar/\Gamma$ when the isolated-pole approximation
is valid.

Definition~\eqref{eq:resolvent-matrix-element} is deliberately phrased in
terms of matrix elements.  The continued resolvent need not be a bounded
operator on the original Hilbert space.  Its poles nevertheless define
finite-rank residues between weighted spaces or between $\PhiSpace$ and its
dual.  This distinction becomes central in the rigged-Hilbert-space
construction of Section~\ref{sec:rigged-hilbert}.

\subsection{Momentum sheets and outgoing boundary conditions}
\label{subsec:siegert}
% BEGIN CURATED SUBJECT INDEX
\index{GNS construction}
\index{partial-wave expansion}
\index{boundary condition!incoming and outgoing}
\index{Riemann sheet}
\index{Gamow--Siegert state}
% END CURATED SUBJECT INDEX

Consider a one-dimensional Hamiltonian
\begin{equation}
  H=-\frac{\hbar^2}{2m}\frac{\rmd^2}{\rmd x^2}+V(x),
  \qquad x\in\real,
  \label{eq:one-dimensional-H}
\end{equation}
where $m>0$ is the particle mass and $V(x)$ is real on the real axis.  Assume
for the moment that $V(x)$ is short ranged.  The energy and momentum are
related by
\begin{equation}
  E=\frac{\hbar^2k^2}{2m} .
  \label{eq:energy-momentum}
\end{equation}
The square root makes the $E$ plane two-sheeted.  On the physical sheet, a
positive-energy scattering solution has real $k>0$.  Analytic continuation
through the cut changes the sign of $k$.

With time dependence $\rme^{-\rmi Et/\hbar}$, a purely outgoing solution obeys
\begin{align}
  \vsub{\psi}{R}(x)
  &\sim A_-\rme^{-\rmi \vsub{k}{R}x},
  &&x\to-\infty,
  \notag\\
  \vsub{\psi}{R}(x)
  &\sim A_+\rme^{+\rmi \vsub{k}{R}x},
  &&x\to+\infty .
  \label{eq:siegert-boundary}
\end{align}
Here $A_-$ and $A_+$ are nonzero amplitudes.  For a decaying resonance,
$\vsub{k}{R}$ lies in the fourth quadrant, so
$\im \vsub{k}{R}<0$.  Both spatial asymptotic forms in
Eq.~\eqref{eq:siegert-boundary} then grow exponentially.  This is not a
pathology that should be removed before defining the pole; it is the spatial
counterpart of temporal decay.  Regularization or analytic deformation is
introduced only after the outgoing condition has selected the physical
solution.

For a spherically symmetric problem, the reduced radial function $u_l(r)$ in
partial wave $l$ satisfies an analogous condition at $r\to\infty$.  The
regular condition at the origin replaces the left outgoing condition.  With
several thresholds, each channel momentum has its own square root, producing
multiple sheets.  A sheet must then be labeled by the signs chosen for all
channel momenta rather than by a single phrase such as ``the second sheet.''

\subsection{Jost functions, the scattering matrix, and pole residues}
\label{subsec:jost}
% BEGIN CURATED SUBJECT INDEX
\index{resolvent}
\index{Fredholm theory}
\index{Jost function}
\index{Wronskian}
\index{scattering matrix}
\index{boundary condition!incoming and outgoing}
\index{resonance!residue}
\index{Gamow--Siegert state}
\index{antiresonance}
\index{Breit--Wigner approximation}
% END CURATED SUBJECT INDEX

We now continue the textbook construction of
Section~\ref{subsec:jost-foundation} into the complex momentum plane.  For a
short-range radial potential, let $f_l(\pm k,r)$ be Jost solutions with
asymptotic behavior
\begin{equation}
  f_l(\pm k,r)\sim\rme^{\pm\rmi kr},
  \qquad r\to\infty .
  \label{eq:jost-asymptotic}
\end{equation}
Let $\varphi_l(k,r)$ be the solution regular at the origin.  Its expansion in
the Jost basis defines the Jost function $F_l(k)$:
\begin{equation}
  \varphi_l(k,r)
  =\frac{1}{2\rmi k}
  \left[F_l(-k)f_l(k,r)-F_l(k)f_l(-k,r)\right] .
  \label{eq:regular-jost-expansion}
\end{equation}
Up to a convention-dependent phase, the partial-wave scattering matrix is
\begin{equation}
  S_l(k)=\frac{F_l(-k)}{F_l(k)} .
  \label{eq:S-jost}
\end{equation}
A zero of $F_l(k)$ removes the incoming component and therefore enforces the
Siegert condition.  The same zero is a pole of $S_l(k)$.  This establishes the
first two arrows in Eq.~\eqref{eq:intro-dictionary}; rigorous versions use
analytic Fredholm theory and weighted resolvents
\cite{Newton:1960,Taylor:1972,Sasakawa:1991,Rakityansky:2022Jost}.

For a real potential, complex conjugation gives the reality relation
\begin{equation}
  F_l(-k^*)=F_l(k)^* .
  \label{eq:jost-reality}
\end{equation}
The conventional classification of simple zeros is summarized in
Table~\ref{tab:jost-zero-classification}.  The terms ``physical'' and
``unphysical'' refer to sheets of the energy surface, not to whether the
corresponding pole has observable consequences.

\begin{table}[t]
  \centering
  \begin{tabular}{lll}
    \toprule
    Position of $k$ & asymptotic behavior & interpretation \\
    \midrule
    $k=+\rmi\kappa$, $\kappa>0$ & $\rme^{-\kappa r}$
      & bound state \\
\addlinespace
    $k=-\rmi\kappa$, $\kappa>0$ & $\rme^{+\kappa r}$
      & virtual or antibound state \\
\addlinespace
    $\re k>0$, $\im k<0$ & outgoing and growing
      & decaying resonance \\
\addlinespace
    $\re k<0$, $\im k<0$ & time-reversed partner
      & antiresonance \\
    \bottomrule
  \end{tabular}
  \caption{Classification of one-channel zeros of $F_l(k)$ for a real
  short-range potential.  Threshold zeros at $k=0$ require a separate
  analysis.}
  \label{tab:jost-zero-classification}
\end{table}

The outgoing radial Green function makes the relation between a Jost zero and
a resolvent pole explicit.  With $r_<:=\min(r,r')$ and
$r_>:=\max(r,r')$, its kernel can be written
\begin{equation}
  G_l^{(+)}(k;r,r')
  =-\frac{2\mu}{\hbar^2}
  \frac{\varphi_l(k,r_<)f_l(k,r_>)}{F_l(k)} .
  \label{eq:radial-green-jost}
\end{equation}
The coefficient follows by integrating the Green equation across $r=r'$ and
using the Wronskian.  If $F_l(k_n)=0$ and $F_l'(k_n)\ne0$, then
\begin{equation}
  \Res_{k=k_n}G_l^{(+)}(k;r,r')
  =-\frac{2\mu}{\hbar^2}
  \frac{\varphi_l(k_n,r_<)f_l(k_n,r_>)}{F_l'(k_n)} .
  \label{eq:radial-green-residue}
\end{equation}
At the zero, $\varphi_l(k_n,r)$ is proportional to $f_l(k_n,r)$, so the
residue factorizes into a right--left dyad.  Passing from a $k$-plane residue
to an energy-plane residue introduces
\begin{equation}
  \frac{\rmd E}{\rmd k}=\frac{\hbar^2k}{\mu} .
  \label{eq:k-to-energy-residue}
\end{equation}
The derivative $F_l'(k_n)$ therefore fixes the pole normalization.  In
Section~\ref{sec:rigged-hilbert} the same normalization is recovered from the
Zel'dovich bilinear integral; the equivalence is a residue identity, not an
accidental agreement of two prescriptions.

The Jost function can also be computed without matching a rapidly growing
solution at a distant boundary.  Write the regular solution as a running
linear combination of incoming and outgoing reference solutions and impose a
variation-of-constants condition that removes second derivatives of the
running coefficients.  The radial Schrodinger equation then becomes a pair
of first-order equations for those coefficients.  Their asymptotic limits are
$F_l(k)$ and $F_l(-k)$.  This formulation extends directly to Jost matrices,
complex angular momentum, Coulomb-modified reference functions, and complex
$k$; it is particularly useful for locating poles numerically
\cite{Rakityansky:2022Jost}.

Near an isolated one-channel pole, unitarity on the real axis and real
analyticity imply the factorization
\begin{equation}
  S(E)=\rme^{2\rmi\vsub{\delta}{bg}(E)}
  \frac{E-\vsub{z}{R}^*}{E-\vsub{z}{R}}
  \vsub{S}{slow}(E) .
  \label{eq:S-pole-factorization}
\end{equation}
The background phase $\vsub{\delta}{bg}$ and the factor
$\vsub{S}{slow}$ vary slowly only when the pole is well isolated from
thresholds and other singularities.  Under that approximation, a cross
section contains the Breit--Wigner denominator
\begin{equation}
  \left|E-\vsub{z}{R}\right|^2
  =(E-E_*)^2+\frac{\Gamma^2}{4} .
  \label{eq:breit-wigner-denominator}
\end{equation}
The residue is as important as the pole position.  It controls excitation
strength, channel branching, and whether interference makes a peak, a dip, or
an asymmetric Fano profile.  In black-hole language the analogous quantities
are excitation factors and residues of the retarded Green function
\cite{Andersson:1995zk,Oshita:2024wgt,Kubota:2025hjk}.

\subsection{Phase shift and time delay}
\label{subsec:time-delay}
% BEGIN CURATED SUBJECT INDEX
\index{phase shift}
\index{Wigner time delay}
\index{scattering!multichannel}
% END CURATED SUBJECT INDEX

For one elastic channel, write $S(E)=\rme^{2\rmi\delta(E)}$, where
$\delta(E)$ is the scattering phase shift.  The Wigner--Smith time delay is
\begin{equation}
  \vsub{\tau}{W}(E)
  =2\hbar\frac{\rmd\delta(E)}{\rmd E} .
  \label{eq:wigner-delay}
\end{equation}
For a unitary multichannel matrix $S(E)$, the corresponding Hermitian matrix is
\begin{equation}
  Q(E)=-\rmi\hbar S^{\dagger}(E)
  \frac{\rmd S(E)}{\rmd E} .
  \label{eq:smith-matrix}
\end{equation}
A narrow isolated resonance contributes a positive Lorentzian to
$\tr Q(E)$, but time delay is not a universal resonance criterion.  Repulsive
backgrounds can yield negative delays, thresholds create non-Lorentzian
singularities, and an exceptionally broad pole may leave no identifiable
maximum.  Pole continuation remains the invariant definition.

\subsection{From poles to time evolution}
\label{subsec:poles-time}
% BEGIN CURATED SUBJECT INDEX
\index{resolvent}
\index{wave packet}
\index{Riemann sheet}
\index{complex scaling}
% END CURATED SUBJECT INDEX

Let $\ket{\chi}$ be a normalized wave packet localized in the interaction
region.  For $t>0$, its survival amplitude can be written as an inverse
Laplace transform of the resolvent:
\begin{equation}
  A_{\chi}(t)
  =\bra{\chi}\rme^{-\rmi Ht/\hbar}\ket{\chi}
  =\frac{1}{2\pi\rmi}
  \int_{\mathcal C}\rmd z\,
  \rme^{-\rmi zt/\hbar}
  \bra{\chi}\Resolvent(z)\ket{\chi} .
  \label{eq:inverse-laplace}
\end{equation}
The contour $\mathcal C$ runs above the spectrum before deformation.  Moving
it onto a nonphysical sheet produces pole terms plus integrals along deformed
cuts.  A simple pole gives Eq.~\eqref{eq:intro-decay}; a threshold at
$E=\vsub{E}{th}$ gives a power-law contribution whose exponent depends on
the threshold density of states.  Thus an operational decomposition is
\begin{equation}
  A_{\chi}(t)
  =\vsub{A}{pole}(t)
  +\vsub{A}{cut}(t)
  +\vsub{A}{arc}(t) .
  \label{eq:time-decomposition}
\end{equation}
The arc term depends on the large-energy deformation and is relevant to the
  short-time expansion.  This contour picture anticipates complex scaling: the
  rotated continuum is the spectral representation of the deformed cut, rather
  than numerical debris to be discarded.
